% page 15 of the facsimile (image p-014 of the scan)
% block 160.0 x 237.7 mm ; left margin 42.0 mm
\pgc{20.320mm}{0.000mm}{
\l{\cel{54}7}
\l{}
\l{}
\l{\sou{adjutanto}.\cel{2}(\sou{armeo}).Oficiro qua servas generalo. - DEFIRS.}
\l{}
\l{\sou{administrar}. (\sou{trans.}) Direktar, surveyar la jero di la aferi}
\l{\cel{3}di societo, di stato, di privato. - DEFIRS.}
\l{}
\l{\sou{admirar}. (\sou{trans, pri, pro})\cel{2}Kontemplar astonate. - EFIS.}
\l{}
\l{\sou{admiralo}. (\sou{armeo navala})\cel{2}Oficiro di qua la grado esas maxima}
\l{\cel{3}en la hierarkio di la navaro militala. - DEFIRS.}
\l{}
\l{\sou{admisar}.\cel{2}(\sou{trans., aden})\cel{2}Aceptar aden loko ulu qua havas la}
\l{\cel{3}qualesi postulata por lo. - EFIS.}
\l{}
\l{\sou{adolecar}.\cel{2}(\sou{netrans}.) Esar en ta evo-periodo qua sequas puer-}
\l{\cel{3}eso e preiras adulteso (t.e. de 15 til 20 ev-yari). - EFIS.}
\l{}
\l{\sou{adoptar}.\cel{2}\sou{(trans.}) I. (\sou{ulu})\cel{2}Prenar legale kom filio. (\sou{Anke}}
\l{\cel{3}\sou{metafore}). II. Igar sua la judiko-maniero o la ago-maniero}
\l{\cel{3}di altru. - DEFIRS.}
\l{}
\l{\sou{adorar}.\cel{2}(\sou{trans})\cel{2}Honorizar quale se lu esus deo. -EFIS.}
\l{}
\l{\sou{adosar}.\cel{2}(\sou{trans., an ulo)}.\cel{2}Apogar per la dorso. - FIS.}
\l{}
\l{\sou{adraganto}. (\sou{bot}.) Gumo qua defluas, forme di paceti vermatra,}
\l{\cel{3}de tragakanto,\cel{2}L.\cel{1}\sou{Astragalus}.\cel{2}DEFIRS.}
\l{}
\l{\sou{adreso}. I. Domicilo di persono o di organismo. - II. Sur-}
\l{\cel{3}skriburo di letro-kuverto, e c., qua indikas ube habitas,}
\l{\cel{3}ube lojas la destinario. - DEFRS.}
\l{}
\l{\sou{adsorbar}.\cel{2}(\sou{trans}.) Koncentrar, sur sua surfaco, substanci}
\l{\cel{3}dissolvita. - DEFI.}
\l{}
\l{\sou{adulta}. Qua atingis la fino di sua kresko. (Dicesas precipue}
\l{\cel{3}pri la homo qua evas plu kam 20 yari). - EFIS.}
\l{}
\l{\sou{adulterar}\cel{1}(\sou{netrans., kun})\cel{2}Violacar sua promiso solena di}
\l{\cel{3}fideleso sexuala a sua spozo. - EFIS.}
\l{}
\l{\sou{adurolo}. (\sou{kemio})\cel{2}Revelilo fotografala, ek klorhidoquinono.}
\l{}
\l{\sou{advento}. (\sou{religio katolika}). La quar semani qui preiras la}
\l{\cel{3}Kristo-nasko - tempo determinita da la Eklezio katolika, dum}
\l{\cel{3}qua la religiani devas preparer su pri la arivo di Yesu-}
\l{\cel{3}Kristo. - DEFIS.}
\l{}
\l{\sou{adventiva}. I. Qua devenas de ago acidenta. - II. (\sou{biol}.)}
\l{\cel{3}\sou{Planto adventiva}\cel{1}: Quan la homo ne semis.}
\l{}
\l{\sou{adverbo}. (\sou{gram}.) Vorto-speco nevariiva qua modifikas la}
\l{\cel{3}senco di adjektivo, di verbo, di adverbo. - DEFIRS.}
}
